\section{Processo di sviluppo adottato} \label{sec:development-process}

\subsection{Metodologia} \label{sec:development-process-methodology}
Il processo di sviluppo che abbiamo scelto di adottare è stato di tipo \textit{Agile} in particolare abbiamo seguito un approccio 
\textit{SCRUM-Like} in modo tale da poter avere un processo di sviluppo iterativo e incrementale che ci permetta di raggiungere obiettivi tangibili
in tempi brevi, sui quali poi costruire le iterazioni successive. Grazie a questo approccio contiamo di agevolare la collaborazione tra i membri
del team, raccogliendo feedback continui sul lavoro svolto e migliorando la qualità del prodotto finale.

Lo sviluppo si svolgerà in sprint della durata di circa una settimana, prima di ogni sprint si svolgerà una riunione di pianificazione
in cui si identificheranno gli obiettivi da raggiungere e le attività da svolgere, dividendo ciascuna di esse in task più piccoli da assegnare 
ai membri del team con il relativo monte ore previsto per lo sviluppo. Durante lo sprint ciascun membro terrà traccia del tempo realmente impiegato
per lo sviluppo di ciascun task, in modo tale da calibrare il lavoro futuro e tenere sotto controllo i tempi di sviluppo. Alla fine
dello sprint si svolgerà una riunione di revisione in cui si presenterà il lavoro svolto da ciascun membro e si raccoglieranno le considerazioni
su cui basare le iterazioni successive. 

\subsubsection{Ruoli e responsabilità} \label{sec:development-process-roles}
Seguendo lo sviluppo SCRUM-Like abbiamo definito i seguenti ruoli e responsabilità all'interno del team:
\begin{itemize}
    \item \textbf{Esperto del dominio} \par
    Francesco Buda è stato designato come esperto del dominio, essendo colui che ha proposto l'idea del progetto.
    \item \textbf{Product Owner} \par
    Tommaso Severi è stato designato come Product Owner, essendo il membro del team con maggiori doti di coordinamento e organizzazione.
    \item \textbf{Sviluppatori} \par
    Ciascun membro del team ha avuto anche il ruolo di sviluppatore.
\end{itemize}

\subsection{Strumenti di sviluppo} \label{sec:development-process-tools}
Per lo sviluppo del progetto sono stati utilizzati i seguenti strumenti:
\begin{itemize}
    \item \textbf{GitFlow} \par
    Per la gestione del codice sorgente e delle versioni del progetto è stato scelto il modello GitFlow, che prevede l'utilizzo di rami
    \texttt{main} e \texttt{develop} rispettivamente per il codice di produzione e sviluppo, oltre che il ramo \texttt{release} per 
    le versioni di rilascio.
    \item \textbf{SBT} \par
    SBT (Scala Build Tool) è stato scelto per la gestione delle dipendenze e la compilazione del progetto Scala.
    \item \textbf{IntelliJ IDEA} \par
    Per lo sviluppo del codice sorgente è stato scelto l'IDE IntelliJ IDEA, che offre un supporto completo per Scala e SBT, oltre a funzionalità 
    utili per un approccio di sviluppo test driven.
\end{itemize}
